% Dark Photon Torsion Model: Physics, Holdom Triviality, and Investigation Findings
% Documents the torsion-trace-as-dark-photon model and the 276-run campaign
% null result, which is the correct physical prediction.
% Uses macros from preamble.tex

\section{Dark Photon Torsion Model}
\label{sec:dark-photon-torsion}

\subsection{Motivation and physics goal}
\label{subsec:dp-motivation}

The torsion dark-photon example
(\texttt{examples/torsion\_dark\_photon/}~and~\texttt{examples/torsion\_dark\_photon\_fv/})
investigates a BSM Lagrangian in which a torsion trace vector is
identified as a Holdom-style dark photon~\cite{Holdom1986}. The goal of
the 276-run HPC campaign was to characterise the region of
$(\alpha, \xi, \delta_m)$ parameter space (or $(m_T^2, \xi, \delta_m)$
in the fundamental-vector formulation) in which a graviton initial
condition produces observable dark-photon conversion via the
combination of Gertsenshtein~\cite{gertsenshtein1962wave} graviton--photon
mixing in a background magnetic field $B_0$ and the
dark-photon kinetic mixing term $\delta_m F \cdot F_t$.

The campaign returned a clean null: every run gave
$P_{\max}(h_5 \to a_1) = \sin^2(\kappa B_0 t / 2) \approx 0.0612$ to
$6.7 \times 10^{-6}$ relative precision, with the torsion (dark-photon)
field amplitudes at numerical round-off. Extensive investigation
(documented below and in issues \#257--\#262) traced this result to
its correct physical origin: \emph{the null is the exact prediction of
this specific Lagrangian}, a manifestation of Holdom's 1986 theorem
that pure kinetic mixing of a photon with a massive dark photon is
unobservable from a pure-photon initial condition. The
276-run campaign result is not a solver bug.

\subsection{Lagrangian}
\label{subsec:dp-lagrangian}

The fundamental-vector form of the Lagrangian (the cleaner of the two
implementations; see~\S\ref{subsec:dp-cdt-vs-fv}) is
\begin{equation}
  \begin{aligned}
    \mathcal{L}
    &= \frac{1}{\kappa^2} R
      - \tfrac{1}{4} F_{\mu\nu} F^{\mu\nu}
      - \tfrac{1}{4} \xi\, F^{(t)}_{\mu\nu} F^{(t)\mu\nu} \\
    &\quad{}
      + \delta_m\, F_{\mu\nu} F^{(t)\mu\nu}
      - \tfrac{1}{2} m_T^2\, t_\mu t^\mu ,
  \end{aligned}
  \label{eq:dp-lagrangian-fv}
\end{equation}
where $A_\mu$ is the Maxwell 4-potential with field strength
$F = \dd A$, $T_\mu$ is the dark-photon 4-vector with field strength
$F^{(t)} = \dd T$, $\xi$ is a dimensionless kinetic coefficient
(canonical normalisation has $\xi = 1$), $\delta_m$ is the
kinetic-mixing parameter, $m_T^2$ is the dark-photon Proca mass, and
$\kappa^2 = 16\pi G$ is the gravitational coupling. (TIDAL uses $\kappa$ as the action prefactor $1/\kappa^2$ on $R$ with raw metric perturbation $g = \eta + h$ — distinct from the Einstein-action $\kappa_\text{E} = 8\pi G$ of some BHL papers; see \cref{sec:kappa-conventions}.) A
$B_0 \hat{x}$-directed background magnetic field is included via
$\bar{A}_\mu = (0, 0, -B_0 z, 0)$.

This Lagrangian matches the conventions used in
\cite{Pospelov2008SecludedU1}, equation~(2):
$\mathcal{L}_{\mathrm{SM}+U(1)'} \supset -\tfrac{1}{4} V_{\mu\nu}^2 +
\tfrac{1}{2}\kappa V_{\mu\nu} F_{\mu\nu}^Y$, with the sign of the
kinetic-mixing cross term \emph{positive}. The sign is conventional
(the transformation $T \to -T$ absorbs a sign flip on $\delta_m$), so
for observable-free dependence on $\delta_m$ the simulation result is
invariant under $\delta_m \to -\delta_m$ --- verified symbolically
in~\S\ref{subsec:dp-holdom-algebra} and numerically.

\subsection{Two formulations: CDT and FV}
\label{subsec:dp-cdt-vs-fv}

The model is implemented in two distinct forms that, after the
solver-side fix discussed in~\S\ref{subsec:dp-solver-fixes},
give numerically identical results at all generic parameter points:

\begin{description}
  \item[CDT (rank-3 trace projection)]
    \texttt{examples/torsion\_dark\_photon/theory.toml}. Uses the
    xPert declaration \texttt{[torsion] perturbation\_name = "t"}
    which linearises the full rank-3 PGT torsion tensor
    $T^a{}_{bc}$ (24 independent components, 12 surviving the
    plane-wave reduction along $\hat{z}$). The Lagrangian couples
    only through the trace $T_\mu = T^\alpha{}_{\mu\alpha}$ and the
    trace-only kinetic $F^{(t)}$ curl. The derived JSON has 18
    dynamical fields with 3 algebraic constraint fields
    ($t_6, t_{13}, t_{20}$) and a rank-deficient mass-matrix block
    on the dynamical torsion components ($t_0, t_{15}, t_{22}$).
  \item[FV (fundamental-vector)]
    \texttt{examples/torsion\_dark\_photon\_fv/theory.toml}. Uses
    \texttt{[[linearization.matter\_perturbations]]} to declare $T$
    as a primitive 4-vector matter field, bypassing the rank-3
    linearisation. The Proca mass is applied to the perturbation
    $t_\mu$ (analogous to the photon-mass convention in
    \texttt{gertsenshtein\_proca}). The derived JSON has 10 clean
    dynamical fields, zero algebraic constraints, and a full-rank
    mass matrix with no hidden gauge DOF.
\end{description}

Both are maintained as a cross-validation pair. At the linearised
level, they describe the same physics: substituting
$(t_0, t_{15}, t_{22}) = (1/3,\, -1/3,\, -1/3)\cdot\psi$ into the CDT
equations reproduces the FV $t_1$ equation exactly (verified
symbolically), with the identification $m_T^2 = -2\alpha$. However,
the CDT JSON encodes the trace direction inside a rank-deficient
24-component tensor structure, and the solver must correctly handle
the rank deficiency to recover the clean physics --- this is the
motivation for the fixes in~\S\ref{subsec:dp-solver-fixes}.

\subsection{Holdom triviality of pure kinetic mixing}
\label{subsec:dp-holdom-algebra}

The central physics finding of this investigation is that the
Lagrangian in~\eqref{eq:dp-lagrangian-fv} produces \emph{zero}
observable photon to dark-photon conversion from a pure-photon
plane-wave initial condition. The mathematical mechanism is an exact
algebraic cancellation in the effective first-order evolution matrix.

Working in the $(a_1, t_1)$ two-field block at wavenumber $k$, the
derived mass and stiffness matrices are
\begin{align}
  M &= \begin{pmatrix} 1 & -2\delta_m \\ +2\delta_m & -\xi \end{pmatrix}
  \label{eq:dp-M} \\
  K &= \begin{pmatrix} -k^2 & +2\delta_m k^2 \\
       -2\delta_m k^2 & m_T^2 + \xi k^2 \end{pmatrix}.
  \label{eq:dp-K}
\end{align}
Note that $M$ is \emph{asymmetric}: the $(a_1, t_1)$ entry
$-2\delta_m$ and the $(t_1, a_1)$ entry $+2\delta_m$ carry opposite
signs. This is a consequence of the sign convention used by the
Wolfram derivation pipeline (moving the $d^2_t$ cross-couplings from
the RHS to the LHS with sign flip) and is the correct representation
of the Lagrangian-derived equations of motion in this basis.

The effective acceleration matrix $E = M^{-1} K$ determines the
first-order evolution $\ddot{\phi} = E \phi$. Computing
$E_{[t_1, a_1]}$ symbolically,
\begin{align}
  E_{[t_1, a_1]}
  &= M^{-1}_{[t_1, a_1]} K_{[a_1, a_1]}
     + M^{-1}_{[t_1, t_1]} K_{[t_1, a_1]} \nonumber\\
  &= \frac{1}{4\delta_m^2 - \xi}\!\left[
       -2\delta_m \cdot (-k^2)
       + 1 \cdot (-2\delta_m k^2) \right] \nonumber\\
  &= \frac{2\delta_m k^2 - 2\delta_m k^2}{4\delta_m^2 - \xi}
   \;\equiv\; 0 .
  \label{eq:dp-cancellation}
\end{align}
The two terms cancel \emph{identically} for any choice of
$\delta_m$, $\xi$, $k$, and $m_T^2$. By contrast, the reverse
coupling is non-zero:
\begin{equation}
  E_{[a_1, t_1]}
  = \frac{2\delta_m m_T^2}{4\delta_m^2 - \xi}
  \label{eq:dp-eff-at}
\end{equation}
which depends linearly on $\delta_m \cdot m_T^2$ --- i.e.\ the
kinetic mixing produces an effective coupling only when combined
with the Proca mass, exactly as predicted by canonical-basis
analysis (see below).

\paragraph{Physical interpretation.}
The cancellation~\eqref{eq:dp-cancellation} says that the photon
$a_1$ is an \emph{exact eigenmode} of the kinetic-mixed system at
the field level. A plane-wave IC pure in $a_1$ is orthogonal to the
dark-photon eigenmode, and no photon $\to$ dark-photon oscillation
occurs. The reverse --- a plane-wave IC pure in $t_1$ --- has
non-trivial projection onto the eigenmode that carries $a_1$
content, and so $a_1$ is excited with substantial amplitude. This
asymmetry has been verified numerically at parameters $(\kappa = 1,
B_0 = 0.01, m_T^2 = 1, \xi = 1, \delta_m = 0.3)$:
\begin{itemize}
  \item $a_1$ plane-wave IC of amplitude $10^{-3}$:
        $t_1$ peak remains at round-off ($\sim 10^{-17}$).
  \item $t_1$ plane-wave IC of amplitude $10^{-3}$:
        $a_1$ peak grows to $\sim 10^{-3}$ (100\,\% amplitude
        transfer).
\end{itemize}

\paragraph{Relation to Holdom 1986.}
This result is the symbolic implementation of Holdom's theorem that
pure kinetic mixing between two $U(1)$ gauge fields is a field
redefinition, and is observable only through coupling to matter
fields (which acquire "millicharges" under the redefinition)
\cite{Holdom1986}. Since our system contains no matter charged
under $U(1)'$, the kinetic mixing is trivially removable by
$T \to T + \delta_m A$ (with appropriate scaling). The solver's
$M^{-1} K$ product performs this field redefinition symbolically via
the matrix inversion, correctly arriving at the decoupled photon
eigenmode.

\paragraph{Relation to Raffelt--Stodolsky conversion.}
The oscillation formula
$P_{\gamma\to\gamma'} = \varepsilon^2 \sin^2(m_{A'}^2 L / 4 k)$
derived in~\cite{An2013StellarConstraints} and elsewhere describes
conversion of a photon \emph{propagating through a medium} (where
the plasma frequency $\omega_p^2$ enters the denominator and breaks
Holdom triviality) or of a photon \emph{produced by a matter source}
(where the source emits in the matter-coupled basis, which after
kinetic-mixing canonicalisation is a superposition of two mass
eigenstates). Neither condition applies to our vacuum
plane-wave IC, which is pumped directly into the $a_1$ field in the
non-canonical basis, i.e.\ exactly the photon eigenmode.

\subsection{Why the 276-run campaign null is correct}
\label{subsec:dp-campaign-interpretation}

The graviton--photon Gertsenshtein coupling in our Lagrangian enters
through the $h_5$ equation as
\begin{equation}
  \ddot{h}_5 \supset \kappa^2 B_0 \partial_x a_1
       - 2 \kappa^2 B_0 \delta_m \partial_x T_{\text{eff}},
\end{equation}
where $T_{\text{eff}}$ is the trace-projected dark photon combination.
The first term is the standard Gertsenshtein source for the photon;
the second term is a direct $B_0 \delta_m$-order coupling between
the graviton and the dark photon.

When the system is pumped with a pure graviton plane-wave IC, the
induced photon $a_1(t)$ lives in the massless photon eigenmode ---
precisely the mode that is decoupled from the dark photon by the
cancellation~\eqref{eq:dp-cancellation}. The direct $h \to T$
coupling (second term above) also passes through the same
$E_{[T, h]} = 0$ cancellation mechanism, giving no dark-photon
excitation even from the direct coupling. Both channels vanish
independently, and the total conversion is exactly zero within
numerical precision.

Therefore the 276-run campaign null is \emph{the correct physical
prediction of this Lagrangian}, not an artefact of the modal
solver. For observable graviton $\to$ dark-photon conversion, the
Lagrangian needs a modification that breaks the Holdom triviality
(see~\S\ref{subsec:dp-what-works}).

\subsection{Solver fixes applied during the investigation}
\label{subsec:dp-solver-fixes}

The investigation uncovered three interlocking modal-solver bugs
that all had to be fixed together before the CDT and FV models
could be shown to give identical results at generic parameters.
They are fixed in commit \texttt{8ddc9cb} on the
\texttt{feature/torsion-sweeps} branch.

\begin{description}
  \item[Hardcoded $M_{\mathrm{mat}}$ diagonal (issue \#258).]
    \texttt{\_build\_evolution\_matrices} previously set
    \texttt{M\_mat[:, fi, fi] = 1.0} unconditionally, ignoring the
    \texttt{kinetic\_coefficient\_symbolic} field of each
    dynamical equation. This caused the off-diagonal structure of
    $M$ to be interpreted inconsistently for equations with non-unit
    kinetic coefficients (e.g.\ $-\kappa^{-2}$ for gravitons,
    $-\xi$ for Proca torsion). Fixed by reading the symbolic
    kinetic coefficient via \texttt{evaluate\_coefficient} and
    installing the resolved value on the diagonal.
  \item[Silent \texttt{np.linalg.eigh} symmetrisation.]
    The rank-deficient Schur elimination path used
    \texttt{eigh} on $M$, assuming it was real symmetric. For the
    kinetic-mixing-derived system, $M$ is asymmetric
    (see~\eqref{eq:dp-M}), and \texttt{eigh} silently symmetrised
    via $(M + M^{\mathsf{T}})/2$, producing a wrong eigenbasis.
    Fixed by replacing with \texttt{np.linalg.svd} plus a
    Schur-elimination step that uses left and right singular
    vectors separately.
  \item[\texttt{\_robust\_eigen\_inv} defensive helper (issue \#259).]
    A fallback \texttt{pinv} path for ill-conditioned eigenvector
    matrices was added earlier to mitigate the IC contamination
    symptom of issue \#258. Once \#258 was root-caused, the
    contamination disappeared and the helper was reverted to
    \texttt{np.linalg.inv}, which is now well-conditioned for
    physical systems.
  \item[\texttt{normalize\_kinetic\_coefficients} CLI band-aid.]
    Three call sites in \texttt{\_simulate.py} and
    \texttt{\_sweep.py} invoked this function to divide each
    equation by its LHS kinetic coefficient, producing a system
    where the hardcoded \texttt{M\_mat[fi,fi] = 1.0} would happen
    to be correct. Removed now that the diagonal is read properly.
\end{description}

\paragraph{Remaining known-issue: rank-deficient critical point
(issue \#260).}
At the specific parameter point $\delta_m = \sqrt{\xi}/2$ (where
$\det M = -\xi + 4\delta_m^2 = 0$ exactly), the $(a_1, t_1)$
block becomes rank-deficient. Both CDT and FV hit the SVD Schur
path at this point and give answers that differ by a factor
$\sim 20$ between the two formulations and from the neighbouring
generic-parameter values. This is a measure-zero parameter
pathology and does not affect the 276-run campaign conclusion,
but it is tracked as an open issue for future cleanup. Sweeps
should avoid tuning to this exact point.

\subsection{What would produce observable conversion}
\label{subsec:dp-what-works}

To break the Holdom triviality in this model, the Lagrangian
would need one of the following modifications, each motivated by a
specific dark-photon physics mechanism.

\begin{description}
  \item[Explicit mass mixing $+\mu^2 A_\mu T^\mu$.]
    An off-diagonal mass-matrix element between $A$ and $T$ makes
    the photon and dark photon genuinely non-eigenmode. The
    symbolic algebra gives
    $E_{[t_1, a_1]} = -\mu^2 / (4\delta_m^2 - \xi)$ which is
    non-zero, and observable conversion proceeds at $O(\mu^2)$.
    This is the "direct mass portal" treatment, distinct from
    kinetic mixing.
  \item[Non-vacuum medium.] In a plasma or medium with effective
    photon mass $\omega_p^2$, the standard Raffelt--Stodolsky
    formula~\cite{An2013StellarConstraints} gives
    $P_{\gamma\to\gamma'} = \varepsilon^2 \sin^2((m_{A'}^2 - \omega_p^2) L / 4k)$
    with $\varepsilon \approx \delta_m$. The photon eigenmode
    acquires an effective mass from the medium, breaking the
    Holdom-trivial decoupling.
  \item[Chern--Simons $\phi F \tilde{F}^{(t)}$.]
    A parity-odd coupling between a pseudoscalar $\phi$ (axion-like)
    and $F\tilde{F}^{(t)}$ produces mixing only in a non-trivial
    background $\bar{F}$, analogous to the
    Raffelt--Stodolsky~\cite{raffelt1988mixing} axion-photon
    mixing.
  \item[Non-minimal gravitational coupling.] A dimension-6 operator
    like $R_{\mu\nu\rho\sigma} F^{\mu\nu} F^{(t)\rho\sigma}$
    couples the dark photon to the curvature tensor, giving a
    graviton-mediated mixing that is non-trivial in the graviton
    background set up by the GW wave itself. This is the
    non-minimal torsion-EM coupling studied in
    \cite{itin2003maxwell} and \cite{rubilar2003torsion}.
  \item[Dark-photon coupling to charged matter.] Adding a coupling
    $\kappa e J_{\mu}^{\mathrm{EM}} T^\mu$ (dark photon couples to
    the electromagnetic current) produces the standard
    ``millicharge'' phenomenology of~\cite{Holdom1986}.
\end{description}

\subsection{Investigation history}
\label{subsec:dp-history}

The investigation that led to this documentation proceeded through
several wrong turns and retractions, all of which are worth noting
as they illustrate the subtlety of the physics:

\begin{enumerate}
  \item Initial (wrong) interpretation: the 276-run null was a
        structural feature of the CDT rank-3 projection, with the
        CDT and FV Lagrangians being "inequivalent at the
        linearised level". \emph{Retracted}: a sympy-based
        symbolic reduction showed they are exactly equivalent
        under the gauge
        $(t_0, t_{15}, t_{22}) = (1/3, -1/3, -1/3)\cdot \psi$.
  \item Intermediate (wrong) claim: a factor-$1/3$ suppression in
        the CDT couplings relative to FV. \emph{Retracted}: a
        projection-normalisation error on my part during the hand
        reduction. With the correct normalisation, the equations
        match term-by-term.
  \item Intermediate (wrong) claim: "FV gives 50/50 channel split"
        at $\delta_m = 0.5$, so the expected physics is
        Raffelt--Stodolsky-like oscillation. \emph{Retracted}: the
        50/50 split was itself a rank-deficient-critical-point
        pathology (issue~\#260); at neighbouring parameters both
        models give pure Gertsenshtein with no dark-photon channel.
  \item Intermediate (wrong) suspicion: the Wolfram-derived
        asymmetric $M_{\mathrm{mat}}$ is a sign bug in the
        derivation pipeline (issue~\#261). \emph{Retracted}: the
        asymmetry is the correct representation of the
        Lagrangian-derived equations in the non-canonical basis.
        It encodes the Holdom-trivial structure precisely: the
        photon eigenmode is pure~$A$ (no~$T$ admixture), which is
        faithfully captured by the $E_{[t, a]} = 0$ cancellation.
  \item Final (correct) picture: the null is real physics. The
        Lagrangian describes a pure kinetic-mixing dark photon
        with Proca mass, which by Holdom's theorem produces no
        observable photon $\to$ dark photon conversion in vacuum
        from a pure-photon IC. To see conversion, the Lagrangian
        needs an explicit mass-mixing term or one of the other
        breakers in~\S\ref{subsec:dp-what-works}.
\end{enumerate}

\subsection{References and GitHub issues}
\label{subsec:dp-references}

Key literature:
\begin{itemize}
  \item B. Holdom, \textit{Two $U(1)$'s and $\varepsilon$ charge
        shifts}, Phys.\ Lett.\ B~\textbf{166} (1986) 196.
        Establishes the field-redefinition triviality of pure
        kinetic mixing.
  \item M. Pospelov, \textit{Secluded $U(1)$ below the weak scale},
        \texttt{arXiv:0811.1030}. Canonical
        Lagrangian~\eqref{eq:dp-lagrangian-fv} with
        positive-sign kinetic-mixing convention.
  \item H. An, M. Pospelov, J. Pradler, \textit{New stellar
        constraints on dark photons},
        \texttt{arXiv:1302.3884}. Eq.~(2.4)--(2.6): vacuum and
        in-medium conversion formulas.
  \item J. Redondo, \textit{Helioscope bounds on hidden sector
        photons}, \texttt{arXiv:0801.1527}.
        Explicit $P_{\gamma\to\gamma'}$ formula with plasma
        frequency.
  \item M. Fabbrichesi, E. Gabrielli, G. Lanfranchi,
        \textit{The Dark Photon}, \texttt{arXiv:2005.01515}.
        Comprehensive review.
  \item A. Caputo, A. Millar, C. O'Hare, E. Vitagliano,
        \textit{Dark photon limits: a cookbook},
        \texttt{arXiv:2105.04565}. Experimental limits and
        oscillation formulae.
  \item G. Raffelt, L. Stodolsky,
        \textit{Mixing of the photon with low-mass particles},
        Phys.\ Rev.\ D~\textbf{37} (1988) 1237.
        Original axion-photon mixing in external $B_0$.
\end{itemize}

Solver-fix and investigation history is tracked in the following
GitHub issues:

\begin{itemize}
  \item \textbf{\#258 (closed)}: hardcoded $M_{\mathrm{mat}}$
        diagonal --- root fix applied in 8ddc9cb.
  \item \textbf{\#259 (closed)}: IC projection $\operatorname{cond}(V)$
        amplification --- root-caused by \#258.
  \item \textbf{\#260 (open)}: CDT vs FV factor-20 discrepancy at
        $\delta_m = \sqrt{\xi}/2$ critical point.
  \item \textbf{\#261 (open)}: asymmetric $M_{\mathrm{mat}}$
        structural investigation --- not a bug, but documentation
        and test coverage needed.
  \item \textbf{\#262 (open)}: this documentation request.
  \item \textbf{\#257 (open)}: rank-deficient eigenbasis --- kept
        open for future genuinely rank-deficient systems beyond
        dark photon.
\end{itemize}

% -- End of section --
